\documentclass[11pt]{article}
% The preamble every project's paper shares: % The preamble every project's paper shares: % The preamble every project's paper shares: \input{../../preamble} from
% projects/<project>/paper/. Boards come from chessboard and xskak, plots from
% pgfplots; every number a paper prints comes from its generated.tex.
\usepackage[LSBC4,T1]{fontenc}
\usepackage[utf8]{inputenc}
\usepackage{lmodern}
\usepackage[margin=1in]{geometry}
\usepackage{microtype}
\usepackage{amsmath}
\usepackage{booktabs}
\usepackage{tabularx}
\usepackage{array}
\usepackage{xcolor}
\usepackage[obeyspaces,spaces]{url}
\usepackage{xskak}
\usepackage{pgfplots}
\pgfplotsset{compat=1.18}
\usepackage{caption}
\usepackage{subcaption}
\usepackage{enumitem}
\usepackage[section]{placeins}
\usepackage[hidelinks]{hyperref}

\setchessboard{boardfontencoding=LSBC4, boardfontfamily=skaknew, blackfieldcolor=black!22,
  whitefieldcolor=white, setfontcolors, showmover=false, label=true, labelfontsize=6pt,
  boardfontsize=14pt}

% A move or line in UCI, and a position in FEN, both breakable at a slash.
\newcommand{\uci}[1]{\texttt{#1}}
\newcommand{\fen}[1]{{\footnotesize\path{#1}}}
\newcolumntype{L}{>{\raggedright\arraybackslash}X}

\setlength{\parskip}{0.4em}
\setlength{\parindent}{0pt}
 from
% projects/<project>/paper/. Boards come from chessboard and xskak, plots from
% pgfplots; every number a paper prints comes from its generated.tex.
\usepackage[LSBC4,T1]{fontenc}
\usepackage[utf8]{inputenc}
\usepackage{lmodern}
\usepackage[margin=1in]{geometry}
\usepackage{microtype}
\usepackage{amsmath}
\usepackage{booktabs}
\usepackage{tabularx}
\usepackage{array}
\usepackage{xcolor}
\usepackage[obeyspaces,spaces]{url}
\usepackage{xskak}
\usepackage{pgfplots}
\pgfplotsset{compat=1.18}
\usepackage{caption}
\usepackage{subcaption}
\usepackage{enumitem}
\usepackage[section]{placeins}
\usepackage[hidelinks]{hyperref}

\setchessboard{boardfontencoding=LSBC4, boardfontfamily=skaknew, blackfieldcolor=black!22,
  whitefieldcolor=white, setfontcolors, showmover=false, label=true, labelfontsize=6pt,
  boardfontsize=14pt}

% A move or line in UCI, and a position in FEN, both breakable at a slash.
\newcommand{\uci}[1]{\texttt{#1}}
\newcommand{\fen}[1]{{\footnotesize\path{#1}}}
\newcolumntype{L}{>{\raggedright\arraybackslash}X}

\setlength{\parskip}{0.4em}
\setlength{\parindent}{0pt}
 from
% projects/<project>/paper/. Boards come from chessboard and xskak, plots from
% pgfplots; every number a paper prints comes from its generated.tex.
\usepackage[LSBC4,T1]{fontenc}
\usepackage[utf8]{inputenc}
\usepackage{lmodern}
\usepackage[margin=1in]{geometry}
\usepackage{microtype}
\usepackage{amsmath}
\usepackage{booktabs}
\usepackage{tabularx}
\usepackage{array}
\usepackage{xcolor}
\usepackage[obeyspaces,spaces]{url}
\usepackage{xskak}
\usepackage{pgfplots}
\pgfplotsset{compat=1.18}
\usepackage{caption}
\usepackage{subcaption}
\usepackage{enumitem}
\usepackage[section]{placeins}
\usepackage[hidelinks]{hyperref}

\setchessboard{boardfontencoding=LSBC4, boardfontfamily=skaknew, blackfieldcolor=black!22,
  whitefieldcolor=white, setfontcolors, showmover=false, label=true, labelfontsize=6pt,
  boardfontsize=14pt}

% A move or line in UCI, and a position in FEN, both breakable at a slash.
\newcommand{\uci}[1]{\texttt{#1}}
\newcommand{\fen}[1]{{\footnotesize\path{#1}}}
\newcolumntype{L}{>{\raggedright\arraybackslash}X}

\setlength{\parskip}{0.4em}
\setlength{\parindent}{0pt}

% Written by projects/his-habits/paper/generate.mjs; do not edit by hand.

\newcommand{\AppCommit}{f4bd63f19b}
\newcommand{\Above}{200}
\newcommand{\Offset}{250}
\newcommand{\WinLow}{10}
\newcommand{\WinHigh}{90}
\newcommand{\ClockFloor}{30}
\newcommand{\CostNodes}{250,000}
\newcommand{\FilterNodes}{120,000}
\newcommand{\PerBand}{2,000}
\newcommand{\PerPlayer}{two}
\newcommand{\Tries}{six}
\newcommand{\LichessFloor}{900}
\newcommand{\MinCount}{30}
\newcommand{\Draws}{2,000}
\newcommand{\KappaLow}{2/3}
\newcommand{\KappaHigh}{3/2}
\newcommand{\Kinds}{22}
\newcommand{\MoveKinds}{ten}
\newcommand{\IdeaKinds}{eleven}
\newcommand{\BandLowest}{1000}
\newcommand{\BandHighest}{1600}
\newcommand{\LichessMonth}{September 2026}
\newcommand{\Games}{647}
\newcommand{\GamesFirst}{324}
\newcommand{\GamesSecond}{323}
\newcommand{\FirstFrom}{July 22}
\newcommand{\FirstTo}{August 23}
\newcommand{\SecondFrom}{August 23}
\newcommand{\SecondTo}{October 15}
\newcommand{\Positions}{9,052}
\newcommand{\PositionsFirst}{4,600}
\newcommand{\PositionsSecond}{4,452}
\newcommand{\SkipOpening}{6,974}
\newcommand{\SkipCheck}{1,614}
\newcommand{\SkipDecided}{2,518}
\newcommand{\SkipClock}{51}
\newcommand{\LichessRead}{121,107}
\newcommand{\LichessKept}{7,962}
\newcommand{\LichessPlayerGames}{14,000}
\newcommand{\LichessWithPositions}{12,256}
\newcommand{\LichessPositions}{23,832}
\newcommand{\HisBandsFirst}{1000 to 1300}
\newcommand{\HisBandsSecond}{1100 to 1400}
\newcommand{\AboveBandsFirst}{1100 to 1500}
\newcommand{\AboveBandsSecond}{1300 to 1600}
\newcommand{\KindTable}{%
\begin{table}[ht]
\centering\small
\begin{tabularx}{\textwidth}{@{}lL@{}}
\toprule
Kind & The move \\
\midrule
Capture & captures \\
Check & gives check \\
Pawn push & moves a pawn without capturing \\
Retreat & steps a knight, bishop, rook or queen back toward its own side, without capturing \\
King move & moves the king, castling aside \\
Queen move & moves the queen \\
Even trade & captures a piece in an exchange the app scores level \\
Worst piece moves & moves the piece the app names his worst \\
Thins his shelter & thins his own king's pawn shelter \\
New threat & leaves him a new capture that wins material, or a new mate in one \\
\midrule
Open file & puts a rook or queen on an open or half-open file \\
Battery & lines up two pieces on a line \\
King pressure & attacks more of the squares around the enemy king, at least three \\
Pawn break & strikes at an enemy pawn chain or centre pawn with a pawn \\
Passed pawn & leaves the moved pawn passed, or promotes it \\
Outpost & settles a knight or bishop on a pawn-guarded square no enemy pawn can attack \\
Seventh rank & lands a rook on the seventh rank against pawns or the king \\
Active king & brings the king forward in an endgame \\
Gives the bishop pair & takes a knight with a bishop and gives up the pair \\
Weakens his pawns & doubles or isolates his own pawns \\
Tactic & forks, pins, skewers, traps a piece or takes one that nothing takes back \\
\midrule
Hangs a piece & leaves a knight, bishop, rook or queen of his standing to be taken, without capturing \\
\bottomrule
\end{tabularx}
\caption{The 22 kinds of move. The first ten are read from the move itself, the next eleven from the insights the app reads for a move, and the last is the control, a known fault at his level.}
\label{tab:kinds}
\end{table}
}
\newcommand{\VerdictTable}{%
\begin{table}[ht]
\centering\footnotesize
\setlength{\tabcolsep}{1.5pt}
\begin{tabular}{@{}lllllr@{}}
\toprule
& \multicolumn{2}{c}{$R$} & \multicolumn{2}{c}{Relative cost (win\%)} & \\
\cmidrule(lr){2-3}\cmidrule(lr){4-5}
Kind & First half & Second half & First half & Second half & Both halves \\
\midrule
Capture & 1.00 {\scriptsize(0.96, 1.03)} & 1.08 {\scriptsize(1.05, 1.12)} & 2.2 {\scriptsize(0.3, 4.1)} & $-$0.9 {\scriptsize($-$2.7, 0.8)} & -- \\
Check & 1.18 {\scriptsize(1.10, 1.26)} & 1.12 {\scriptsize(1.03, 1.22)} & $-$2.0 {\scriptsize($-$5.0, 1.0)} & $-$1.8 {\scriptsize($-$4.6, 1.3)} & more often, costs no more \\
Pawn push & 1.11 {\scriptsize(1.04, 1.18)} & 0.96 {\scriptsize(0.89, 1.04)} & 0.4 {\scriptsize($-$1.0, 1.9)} & 0.0 {\scriptsize($-$1.2, 1.2)} & -- \\
Retreat & 0.86 {\scriptsize(0.80, 0.92)} & 0.88 {\scriptsize(0.81, 0.95)} & 0.5 {\scriptsize($-$1.4, 2.5)} & 0.5 {\scriptsize($-$1.6, 2.6)} & less often, costs no more \\
King move & 0.91 {\scriptsize(0.78, 1.02)} & 0.96 {\scriptsize(0.83, 1.11)} & 0.2 {\scriptsize($-$2.0, 2.8)} & 0.7 {\scriptsize($-$2.6, 4.8)} & -- \\
Queen move & 0.99 {\scriptsize(0.94, 1.03)} & 0.98 {\scriptsize(0.92, 1.03)} & 2.0 {\scriptsize(0.0, 4.1)} & 0.2 {\scriptsize($-$1.7, 2.2)} & -- \\
Even trade & 1.08 {\scriptsize(0.97, 1.20)} & 1.18 {\scriptsize(1.05, 1.32)} & $-$0.8 {\scriptsize($-$5.5, 4.2)} & $-$3.7 {\scriptsize($-$6.9, $-$0.6)} & -- \\
Worst piece moves & 0.98 {\scriptsize(0.92, 1.05)} & 1.06 {\scriptsize(0.98, 1.13)} & $-$0.2 {\scriptsize($-$2.4, 2.1)} & 0.2 {\scriptsize($-$1.3, 1.8)} & -- \\
Thins his shelter & 1.11 {\scriptsize(0.97, 1.28)} & 1.02 {\scriptsize(0.87, 1.18)} & 2.7 {\scriptsize($-$0.0, 5.5)} & 0.9 {\scriptsize($-$1.0, 2.8)} & -- \\
New threat & 0.99 {\scriptsize(0.95, 1.02)} & 0.98 {\scriptsize(0.94, 1.02)} & $-$0.0 {\scriptsize($-$1.4, 1.3)} & 0.3 {\scriptsize($-$1.0, 1.6)} & -- \\
\midrule
Open file & 1.05 {\scriptsize(0.98, 1.13)} & 0.95 {\scriptsize(0.86, 1.04)} & $-$1.8 {\scriptsize($-$3.3, $-$0.1)} & $-$2.2 {\scriptsize($-$4.0, $-$0.5)} & -- \\
Battery & 1.04 {\scriptsize(0.92, 1.18)} & 0.91 {\scriptsize(0.79, 1.06)} & 0.1 {\scriptsize($-$2.7, 3.2)} & 1.3 {\scriptsize($-$1.3, 4.2)} & -- \\
King pressure & 0.95 {\scriptsize(0.85, 1.08)} & 1.08 {\scriptsize(0.94, 1.25)} & $-$1.6 {\scriptsize($-$5.0, 1.5)} & $-$0.4 {\scriptsize($-$4.4, 3.4)} & -- \\
Pawn break & 1.60 {\scriptsize(1.30, 1.97)} & 1.51 {\scriptsize(1.21, 1.94)} & 1.7 {\scriptsize($-$0.6, 4.1)} & 0.7 {\scriptsize($-$1.1, 2.6)} & more often, costs no more \\
Passed pawn & 1.15 {\scriptsize(0.97, 1.33)} & 1.01 {\scriptsize(0.84, 1.20)} & $-$0.7 {\scriptsize($-$5.5, 3.9)} & $-$1.9 {\scriptsize($-$5.4, 1.3)} & -- \\
Outpost & 0.75 {\scriptsize(0.51, 1.02)} & 0.65 {\scriptsize(0.45, 1.05)} & 3.6 {\scriptsize($-$0.5, 8.0)} & $-$0.9 {\scriptsize($-$3.1, 1.7)} & -- \\
Seventh rank & 1.16 {\scriptsize(0.93, 1.46)} & 1.23 {\scriptsize(0.90, 1.59)} & $-$4.7 {\scriptsize($-$8.3, $-$1.8)} & 0.6 {\scriptsize($-$4.7, 7.2)} & -- \\
Active king & 1.01 {\scriptsize(0.80, 1.35)} & 1.00 {\scriptsize(0.71, 1.46)} & 1.0 {\scriptsize($-$6.1, 11.5)} & $-$3.8 {\scriptsize($-$6.4, $-$1.8)} & -- \\
Gives the bishop pair & 1.04 {\scriptsize(0.78, 1.36)} & 1.11 {\scriptsize(0.87, 1.49)} & 1.0 {\scriptsize($-$4.4, 7.8)} & $-$1.4 {\scriptsize($-$6.5, 3.9)} & too few ($E=29.9$) \\
Weakens his pawns & 1.13 {\scriptsize(1.02, 1.25)} & 1.28 {\scriptsize(1.16, 1.43)} & 0.5 {\scriptsize($-$1.8, 3.0)} & 0.9 {\scriptsize($-$1.8, 3.7)} & more often, costs no more \\
Tactic & 1.03 {\scriptsize(0.97, 1.08)} & 0.94 {\scriptsize(0.89, 1.00)} & $-$2.6 {\scriptsize($-$5.9, 0.4)} & 0.2 {\scriptsize($-$2.1, 2.7)} & -- \\
\midrule
Hangs a piece & 0.90 {\scriptsize(0.73, 1.09)} & 0.93 {\scriptsize(0.70, 1.20)} & 13.8 {\scriptsize(9.3, 18.6)} & 14.1 {\scriptsize(9.1, 19.4)} & -- \\
\bottomrule
\end{tabular}
\caption{Each kind of move against Maia-2 at the band above. $R$ is his count over the expected count corrected by Maia-2's bias on Lichess, with its 95\% interval. The relative cost is the mean win\% his instances give up against the band above's most likely move, minus the same over his other differences: for a kind he plays more often on both halves, or with no such reading where the first half's $R$ is at least 1, the positions where he plays it and the band above would not; otherwise the positions where the band above would and he does not. For the control, the direction it tests: he leaves a piece hanging and the band above would not. The last column is what the rule would have read had the controls passed: the kinds whose $R$ lies on one side of 1 on both halves, and their cost. The known fault failed, so none of it is a verdict.}
\label{tab:verdicts}
\end{table}
}
\newcommand{\LevelTable}{%
\begin{table}[ht]
\centering\small
\setlength{\tabcolsep}{3pt}
\begin{tabular}{@{}lllll@{}}
\toprule
& \multicolumn{2}{c}{$R$ against his own band} & \multicolumn{2}{c}{Maia-2's correction} \\
\cmidrule(lr){2-3}\cmidrule(lr){4-5}
Kind & First half & Second half & First half & Second half \\
\midrule
Capture & 1.02 {\scriptsize(0.99, 1.05)} & 1.06 {\scriptsize(1.02, 1.10)} & 0.99 & 0.98 \\
Check & 1.19 {\scriptsize(1.11, 1.27)} & 1.06 {\scriptsize(0.97, 1.15)} & 1.02 & 1.00 \\
Pawn push & 1.05 {\scriptsize(0.98, 1.12)} & 0.90 {\scriptsize(0.82, 0.97)} & 1.05 & 1.03 \\
Retreat & 0.92 {\scriptsize(0.85, 0.99)} & 0.93 {\scriptsize(0.85, 1.02)} & 0.91 & 0.91 \\
King move & 1.01 {\scriptsize(0.86, 1.16)} & 1.02 {\scriptsize(0.88, 1.17)} & 1.00 & 1.06 \\
Queen move & 1.01 {\scriptsize(0.96, 1.06)} & 1.00 {\scriptsize(0.95, 1.06)} & 0.97 & 0.97 \\
Even trade & 1.16 {\scriptsize(1.03, 1.29)} & 1.13 {\scriptsize(1.01, 1.27)} & 1.01 & 0.99 \\
Worst piece moves & 0.97 {\scriptsize(0.90, 1.04)} & 1.03 {\scriptsize(0.95, 1.11)} & 0.98 & 0.96 \\
Thins his shelter & 1.08 {\scriptsize(0.94, 1.24)} & 0.93 {\scriptsize(0.80, 1.09)} & 1.01 & 1.00 \\
New threat & 0.99 {\scriptsize(0.96, 1.03)} & 0.99 {\scriptsize(0.95, 1.04)} & 1.02 & 1.02 \\
\midrule
Open file & 1.03 {\scriptsize(0.96, 1.11)} & 0.99 {\scriptsize(0.91, 1.08)} & 1.00 & 1.02 \\
Battery & 1.06 {\scriptsize(0.94, 1.22)} & 0.98 {\scriptsize(0.84, 1.14)} & 0.93 & 0.95 \\
King pressure & 0.90 {\scriptsize(0.80, 1.01)} & 1.11 {\scriptsize(0.97, 1.29)} & 0.95 & 0.97 \\
Pawn break & 1.38 {\scriptsize(1.14, 1.67)} & 1.45 {\scriptsize(1.14, 1.86)} & 0.96 & 0.97 \\
Passed pawn & 1.16 {\scriptsize(0.99, 1.35)} & 1.24 {\scriptsize(1.02, 1.51)} & 0.97 & 1.18 \\
Outpost & 0.79 {\scriptsize(0.55, 1.13)} & 0.82 {\scriptsize(0.55, 1.12)} & 0.99 & 1.22 \\
Seventh rank & 1.13 {\scriptsize(0.93, 1.44)} & 1.15 {\scriptsize(0.88, 1.52)} & 1.01 & 0.92 \\
Active king & 1.16 {\scriptsize(0.94, 1.50)} & 0.85 {\scriptsize(0.58, 1.16)} & 1.17 & 1.06 \\
Gives the bishop pair & 1.07 {\scriptsize(0.81, 1.44)} & 1.13 {\scriptsize(0.90, 1.44)} & 1.08 & 1.10 \\
Weakens his pawns & 1.18 {\scriptsize(1.07, 1.30)} & 1.25 {\scriptsize(1.13, 1.38)} & 0.99 & 0.96 \\
Tactic & 1.05 {\scriptsize(0.99, 1.11)} & 1.01 {\scriptsize(0.95, 1.07)} & 0.99 & 1.06 \\
\midrule
Hangs a piece & 0.77 {\scriptsize(0.63, 0.93)} & 0.67 {\scriptsize(0.52, 0.86)} & 1.13 & 0.97 \\
\bottomrule
\end{tabular}
\caption{Each kind against Maia-2 at his own band, corrected the same way, and the correction itself: the corrected expected count at the band above over Maia-2's own. A correction above 1 means players at those bands play the kind more often than Maia-2 predicts.}
\label{tab:level}
\end{table}
}
\newcommand{\CostTable}{%
\begin{table}[ht]
\centering\small
\setlength{\tabcolsep}{4pt}
\begin{tabular}{@{}lrrrrrrrr@{}}
\toprule
& \multicolumn{4}{c}{He plays it, the band above would not} & \multicolumn{4}{c}{The band above would, he does not} \\
\cmidrule(lr){2-5}\cmidrule(lr){6-9}
& \multicolumn{2}{c}{First half} & \multicolumn{2}{c}{Second half} & \multicolumn{2}{c}{First half} & \multicolumn{2}{c}{Second half} \\
\cmidrule(lr){2-3}\cmidrule(lr){4-5}\cmidrule(lr){6-7}\cmidrule(lr){8-9}
Kind & A game & $D$ & A game & $D$ & A game & $D$ & A game & $D$ \\
\midrule
Capture & 0.40 & 3.7 & 0.51 & 4.9 & 1.49 & 5.1 & 0.95 & 2.3 \\
Check & 0.48 & 1.6 & 0.34 & 1.5 & 0.50 & 4.3 & 0.44 & 4.4 \\
Pawn push & 1.44 & 4.0 & 1.12 & 3.2 & 0.85 & 3.2 & 0.84 & 3.4 \\
Retreat & 0.62 & 3.4 & 0.60 & 2.9 & 0.73 & 4.0 & 0.86 & 3.7 \\
King move & 0.27 & 6.6 & 0.24 & 3.6 & 0.24 & 3.6 & 0.14 & 3.9 \\
Queen move & 0.71 & 4.7 & 0.79 & 4.5 & 0.67 & 5.2 & 0.78 & 3.3 \\
Even trade & 0.16 & 2.6 & 0.20 & $-$0.3 & 0.42 & 5.3 & 0.33 & 0.9 \\
Worst piece moves & 0.70 & 3.1 & 0.71 & 2.8 & 0.79 & 3.4 & 0.80 & 3.5 \\
Thins his shelter & 0.41 & 6.0 & 0.28 & 4.0 & 0.26 & 4.3 & 0.19 & 4.1 \\
New threat & 1.28 & 3.6 & 1.13 & 3.9 & 1.95 & 3.3 & 1.74 & 2.9 \\
\midrule
Open file & 0.86 & 1.9 & 0.94 & 1.0 & 0.66 & 4.4 & 0.90 & 5.1 \\
Battery & 0.30 & 3.7 & 0.37 & 4.3 & 0.36 & 2.3 & 0.54 & 5.0 \\
King pressure & 0.21 & 4.2 & 0.23 & 2.9 & 0.35 & 2.0 & 0.25 & 2.8 \\
Pawn break & 0.30 & 5.2 & 0.28 & 3.9 & 0.12 & $-$0.2 & 0.14 & 2.7 \\
Passed pawn & 0.18 & 2.8 & 0.16 & 1.3 & 0.24 & 5.9 & 0.14 & 5.3 \\
Outpost & 0.02 & $-$7.0 & 0.02 & 0.7 & 0.08 & 7.1 & 0.06 & 2.3 \\
Seventh rank & 0.07 & $-$1.1 & 0.08 & 3.8 & 0.07 & 2.7 & 0.14 & 3.6 \\
Active king & 0.09 & 4.5 & 0.03 & $-$0.6 & 0.08 & 3.8 & 0.04 & 4.8 \\
Gives the bishop pair & 0.05 & 4.5 & 0.07 & 1.8 & 0.06 & 3.0 & 0.05 & 5.5 \\
Weakens his pawns & 0.19 & 3.9 & 0.17 & 4.1 & 0.29 & 5.2 & 0.14 & 4.8 \\
Tactic & 0.32 & 0.8 & 0.29 & 3.5 & 0.42 & 7.4 & 0.40 & 3.2 \\
\midrule
Hangs a piece & 0.31 & 16.9 & 0.22 & 16.9 & 0.07 & $-$7.2 & 0.08 & $-$4.5 \\
\bottomrule
\end{tabular}
\caption{The instances of each kind in each direction: how many a game, and their mean $D$, the win\% his move gives up against the band above's most likely move. Over all the positions where his move is not that move, his move gives up 3.5 win\% on average in the first half and 3.2 in the second; a kind's reference is those positions less the ones where either move is of the kind.}
\label{tab:cost}
\end{table}
}
\newcommand{\KappaTable}{%
\begin{table}[ht]
\centering\small
\setlength{\tabcolsep}{3pt}
\begin{tabular}{@{}lrrrrrrr@{}}
\toprule
Kind & 1000\textsuperscript{a} & 1100 & 1200 & 1300 & 1400 & 1500 & 1600 \\
\midrule
Capture & 0.96 & 0.97 & 1.03 & 0.99 & 0.98 & 0.98 & 0.99 \\
Check & 1.06 & 0.99 & 1.05 & 1.05 & 0.98 & 1.01 & 1.04 \\
Pawn push & 1.05 & 1.04 & 1.06 & 1.06 & 1.04 & 1.03 & 1.08 \\
Retreat & 0.85 & 0.96 & 0.86 & 0.89 & 0.93 & 0.90 & 0.96 \\
King move & 0.92 & 0.92 & 1.00 & 1.02 & 1.06 & 1.06 & 0.90 \\
Queen move & 0.98 & 0.97 & 0.95 & 0.96 & 1.00 & 0.96 & 0.98 \\
Even trade & 0.99 & 0.96 & 1.05 & 1.00 & 1.02 & 0.98 & 1.03 \\
Worst piece moves & 1.07 & 0.96 & 0.97 & 0.99 & 0.99 & 0.95 & 0.96 \\
Thins his shelter & 1.05 & 1.06 & 0.90 & 1.10 & 0.95 & 1.00 & 1.02 \\
New threat & 1.02 & 1.04 & 0.98 & 1.01 & 1.03 & 1.02 & 1.02 \\
\midrule
Open file & 1.12 & 0.98 & 1.01 & 1.01 & 0.99 & 1.02 & 0.99 \\
Battery & 1.03 & 0.87 & 0.91 & 0.91 & 1.01 & 0.93 & 0.98 \\
King pressure & 1.03 & 1.00 & 1.03 & 0.91 & 0.92 & 0.98 & 0.99 \\
Pawn break & 1.17 & 0.83 & 1.14 & 0.90 & 1.06 & 0.95 & 1.10 \\
Passed pawn & 0.93 & 0.94 & 0.93 & 0.95 & 1.05 & 1.21 & 1.01 \\
Outpost & 0.99\textsuperscript{b} & 0.99\textsuperscript{b} & 0.99\textsuperscript{b} & 0.99\textsuperscript{b} & 0.99\textsuperscript{b} & 1.25 & 0.90 \\
Seventh rank & 1.14 & 1.01\textsuperscript{b} & 1.01\textsuperscript{b} & 0.99 & 1.07 & 0.88 & 1.05 \\
Active king & 1.06\textsuperscript{b} & 1.06\textsuperscript{b} & 1.18 & 1.27 & 1.04 & 1.06\textsuperscript{b} & 1.02 \\
Gives the bishop pair & 1.05\textsuperscript{b} & 1.05\textsuperscript{b} & 1.15 & 1.05\textsuperscript{b} & 1.09 & 1.11 & 0.90 \\
Weakens his pawns & 0.88 & 0.96 & 1.04 & 0.99 & 0.95 & 0.97 & 0.98 \\
Tactic & 1.00 & 0.96 & 0.99 & 0.98 & 1.03 & 1.06 & 1.01 \\
\midrule
Hangs a piece & 1.17 & 1.07 & 1.15 & 1.22 & 1.04 & 0.95 & 1.01 \\
\bottomrule
\end{tabular}
\caption{$\kappa$, Maia-2's bias at each band on Lichess: the moves of a kind players at the band played over those Maia-2 at the band expected. \textsuperscript{a}Maia-2's band below 1100, here from 900. \textsuperscript{b}Maia-2's expected count at that band is under 30, so the bias pooled over the bands stands in.}
\label{tab:kappa}
\end{table}
}
\newcommand{\PhaseTable}{%
\begin{table}[ht]
\centering\small
\begin{tabular}{@{}lll@{}}
\toprule
Kind & Middlegame & Endgame \\
\midrule
Capture & 1.05 {\scriptsize(1.02, 1.08)} & 1.00 {\scriptsize(0.96, 1.05)} \\
Check & 1.18 {\scriptsize(1.10, 1.27)} & 1.12 {\scriptsize(1.05, 1.20)} \\
Pawn push & 1.00 {\scriptsize(0.94, 1.06)} & 1.13 {\scriptsize(1.04, 1.22)} \\
Retreat & 0.83 {\scriptsize(0.78, 0.89)} & 0.94 {\scriptsize(0.87, 1.02)} \\
King move & 0.70 {\scriptsize(0.55, 0.86)} & 1.00 {\scriptsize(0.90, 1.10)} \\
Queen move & 0.98 {\scriptsize(0.94, 1.02)} & 0.98 {\scriptsize(0.93, 1.03)} \\
Even trade & 1.15 {\scriptsize(1.05, 1.24)} & 1.02 {\scriptsize(0.80, 1.25)} \\
Worst piece moves & 1.04 {\scriptsize(0.97, 1.12)} & 1.00 {\scriptsize(0.95, 1.06)} \\
Thins his shelter & 1.04 {\scriptsize(0.92, 1.17)} & 1.12 {\scriptsize(0.97, 1.29)} \\
New threat & 0.99 {\scriptsize(0.96, 1.02)} & 0.96 {\scriptsize(0.91, 1.01)} \\
\midrule
Open file & 1.01 {\scriptsize(0.94, 1.08)} & 0.98 {\scriptsize(0.91, 1.05)} \\
Battery & 0.95 {\scriptsize(0.85, 1.05)} & 1.03 {\scriptsize(0.89, 1.18)} \\
King pressure & 1.07 {\scriptsize(0.95, 1.20)} & 0.96 {\scriptsize(0.86, 1.09)} \\
Pawn break & 1.56 {\scriptsize(1.33, 1.84)} & 1.53 {\scriptsize(1.11, 2.09)} \\
Passed pawn & 0.97 {\scriptsize(0.83, 1.14)} & 1.16 {\scriptsize(1.01, 1.33)} \\
Outpost & 0.74 {\scriptsize(0.56, 1.03)} & 0.58 {\scriptsize(0.30, 1.01)} \\
Seventh rank & 1.07 {\scriptsize(0.81, 1.33)} & 1.30 {\scriptsize(1.04, 1.59)} \\
Active king & -- & 1.00 {\scriptsize(0.84, 1.28)} \\
Gives the bishop pair & 1.08 {\scriptsize(0.88, 1.35)} & 1.41 \\
Weakens his pawns & 1.21 {\scriptsize(1.12, 1.31)} & 1.15 {\scriptsize(1.02, 1.31)} \\
Tactic & 0.98 {\scriptsize(0.93, 1.02)} & 1.00 {\scriptsize(0.93, 1.06)} \\
\midrule
Hangs a piece & 0.92 {\scriptsize(0.76, 1.09)} & 0.88 {\scriptsize(0.63, 1.18)} \\
\bottomrule
\end{tabular}
\caption{$R$ against the band above by the divider's phase of the position, both halves together.}
\label{tab:phase}
\end{table}
}
\newcommand{\HangsRFirst}{0.90}
\newcommand{\HangsRSecond}{0.93}
\newcommand{\HangsRFirstCi}{0.73 to 1.09}
\newcommand{\HangsRSecondCi}{0.70 to 1.20}
\newcommand{\HangsRhisFirst}{0.77}
\newcommand{\HangsRhisSecond}{0.67}
\newcommand{\HangsRhisFirstCi}{0.63 to 0.93}
\newcommand{\HangsRhisSecondCi}{0.52 to 0.86}
\newcommand{\HangsNFirst}{111}
\newcommand{\HangsNSecond}{82}
\newcommand{\HangsOverDFirst}{16.9}
\newcommand{\HangsOverDSecond}{16.9}
\newcommand{\HangsRelFirst}{13.8}
\newcommand{\HangsRelSecond}{14.1}
\newcommand{\HangsRelFirstCi}{9.3 to 18.6}
\newcommand{\HangsRelSecondCi}{9.1 to 19.4}
\newcommand{\StepFirst}{1.17}
\newcommand{\StepSecond}{1.37}
\newcommand{\SpreadFirst}{19}
\newcommand{\SpreadSecond}{27}
\newcommand{\NullMisses}{three}
\newcommand{\NullIntervals}{44}
\newcommand{\TopAboveFirst}{51}
\newcommand{\TopAboveSecond}{52}
\newcommand{\TopHisFirst}{51}
\newcommand{\TopHisSecond}{53}
\newcommand{\DiffShareFirst}{49}
\newcommand{\DiffShareSecond}{48}
\newcommand{\DiffDFirst}{3.5}
\newcommand{\DiffDSecond}{3.2}
\newcommand{\OneHalf}{four}
\newcommand{\BreakRFirst}{1.60}
\newcommand{\BreakRSecond}{1.51}
\newcommand{\BreakRFirstCi}{1.30 to 1.97}
\newcommand{\BreakRSecondCi}{1.21 to 1.94}
\newcommand{\BreakRhisFirst}{1.38}
\newcommand{\BreakRhisSecond}{1.45}
\newcommand{\BreakRhisFirstCi}{1.14 to 1.67}
\newcommand{\BreakRhisSecondCi}{1.14 to 1.86}
\newcommand{\RetreatRFirst}{0.86}
\newcommand{\RetreatRSecond}{0.88}
\newcommand{\RetreatRFirstCi}{0.80 to 0.92}
\newcommand{\RetreatRSecondCi}{0.81 to 0.95}
\newcommand{\RetreatRhisFirst}{0.92}
\newcommand{\RetreatRhisSecond}{0.93}
\newcommand{\RetreatRhisFirstCi}{0.85 to 0.99}
\newcommand{\RetreatRhisSecondCi}{0.85 to 1.02}
\newcommand{\CheckRFirst}{1.18}
\newcommand{\CheckRSecond}{1.12}
\newcommand{\CheckRFirstCi}{1.10 to 1.26}
\newcommand{\CheckRSecondCi}{1.03 to 1.22}
\newcommand{\CheckRhisFirst}{1.19}
\newcommand{\CheckRhisSecond}{1.06}
\newcommand{\CheckRhisFirstCi}{1.11 to 1.27}
\newcommand{\CheckRhisSecondCi}{0.97 to 1.15}
\newcommand{\WeakRFirst}{1.13}
\newcommand{\WeakRSecond}{1.28}
\newcommand{\WeakRFirstCi}{1.02 to 1.25}
\newcommand{\WeakRSecondCi}{1.16 to 1.43}
\newcommand{\WeakRhisFirst}{1.18}
\newcommand{\WeakRhisSecond}{1.25}
\newcommand{\WeakRhisFirstCi}{1.07 to 1.30}
\newcommand{\WeakRhisSecondCi}{1.13 to 1.38}
\newcommand{\TradeRFirst}{1.08}
\newcommand{\TradeRSecond}{1.18}
\newcommand{\TradeRFirstCi}{0.97 to 1.20}
\newcommand{\TradeRSecondCi}{1.05 to 1.32}
\newcommand{\TradeRhisFirst}{1.16}
\newcommand{\TradeRhisSecond}{1.13}
\newcommand{\TradeRhisFirstCi}{1.03 to 1.29}
\newcommand{\TradeRhisSecondCi}{1.01 to 1.27}
\newcommand{\BreakDFirst}{5.2}
\newcommand{\BreakDSecond}{3.9}
\newcommand{\BreakPerGameFirst}{0.30}
\newcommand{\BreakPerGameSecond}{0.28}
\newcommand{\RetreatDFirst}{4.0}
\newcommand{\RetreatDSecond}{3.7}
\newcommand{\RetreatPerGameFirst}{0.73}
\newcommand{\RetreatPerGameSecond}{0.86}
\newcommand{\CheckDFirst}{1.6}
\newcommand{\CheckDSecond}{1.5}
\newcommand{\RetreatMid}{0.83}
\newcommand{\RetreatMidCi}{0.78 to 0.89}
\newcommand{\RetreatEnd}{0.94}
\newcommand{\RetreatEndCi}{0.87 to 1.02}
\newcommand{\BreakMid}{1.56}
\newcommand{\BreakEnd}{1.53}
\newcommand{\KingMid}{0.70}
\newcommand{\KingMidCi}{0.55 to 0.86}
\newcommand{\KappaMin}{0.83}
\newcommand{\KappaMax}{1.27}
\newcommand{\KappaPooled}{13}
\newcommand{\KappaCells}{154}
\newcommand{\CorrMin}{0.91}
\newcommand{\CorrMax}{1.22}
\newcommand{\LiHangsLow}{3.9}
\newcommand{\LiHangsHigh}{1.9}
\newcommand{\LiRetreatLow}{9.3}
\newcommand{\LiRetreatHigh}{12.2}
\newcommand{\LiCheckLow}{10.4}
\newcommand{\LiCheckHigh}{7.8}
\newcommand{\LiBreakMin}{1.4}
\newcommand{\LiBreakMax}{2.2}
\newcommand{\HisBreakFirst}{2.9}
\newcommand{\HisBreakSecond}{2.7}
\newcommand{\MaiaBreakFirst}{1.9}
\newcommand{\MaiaBreakSecond}{1.8}
\newcommand{\LikeFirst}{0.9}
\newcommand{\LikeSecond}{2.1}
\newcommand{\HisHangsFirst}{2.4}
\newcommand{\HisHangsSecond}{1.8}
\newcommand{\MaiaHangsHisFirst}{2.8}
\newcommand{\MaiaHangsHisSecond}{2.3}
\newcommand{\MaiaHangsAboveFirst}{2.4}
\newcommand{\MaiaHangsAboveSecond}{2.0}
\newcommand{\TrendTable}{%
\begin{table}[ht]
\centering\small
\setlength{\tabcolsep}{4pt}
\begin{tabular}{@{}lrrrrrrrrr@{}}
\toprule
& \multicolumn{7}{c}{Lichess players at each band} & \multicolumn{2}{c}{Him} \\
\cmidrule(lr){2-8}\cmidrule(lr){9-10}
Kind & 1000\textsuperscript{a} & 1100 & 1200 & 1300 & 1400 & 1500 & 1600 & First & Second \\
\midrule
Pawn break & 2.1 & 1.4 & 1.9 & 1.7 & 1.8 & 1.8 & 2.2 & 2.9 & 2.7 \\
Retreat & 9.3 & 11.1 & 9.8 & 10.2 & 11.4 & 11.5 & 12.2 & 10.5 & 12.1 \\
Check & 10.4 & 8.8 & 9.2 & 9.3 & 8.5 & 7.9 & 7.8 & 12.3 & 10.4 \\
Weakens his pawns & 4.7 & 5.1 & 5.0 & 4.7 & 4.6 & 4.6 & 4.1 & 5.3 & 5.3 \\
Even trade & 3.8 & 4.8 & 4.9 & 5.1 & 5.4 & 5.2 & 5.4 & 5.0 & 5.5 \\
Hangs a piece & 3.9 & 2.8 & 3.0 & 3.0 & 2.0 & 1.9 & 1.9 & 2.4 & 1.8 \\
\bottomrule
\end{tabular}
\caption{Added after the run. The share of positions, in \%, in which players play each kind: Lichess players at each band in their own positions, and him in his. The positions differ, so the columns compare the bands with each other, not with him. \textsuperscript{a}Maia-2's band below 1100, here from 900.}
\label{tab:trend}
\end{table}
}
\newcommand{\OneStepGames}{62}
\newcommand{\GamesWithPositions}{601}
\newcommand{\EaFen}{\fen{5rk1/4qppp/4p3/3p3b/3P4/3QP3/Pr3PPP/R3R1K1 w - - 0 22}}
\newcommand{\EaMover}{White}
\newcommand{\EaPhase}{middlegame}
\newcommand{\EaPlayed}{22.~e4}
\newcommand{\EaTop}{22.~a4}
\newcommand{\EaPlayedUci}{\uci{e3e4}}
\newcommand{\EaTopUci}{\uci{a2a4}}
\newcommand{\EaLineP}{22.~e4 \symqueen{}f6 23.~\symqueen{}e3 dxe4}
\newcommand{\EaLineT}{22.~a4 \symqueen{}f6 23.~\symqueen{}f1 \symbishop{}g6}
\newcommand{\EaWpBefore}{22}
\newcommand{\EaWpP}{21}
\newcommand{\EaWpT}{24}
\newcommand{\EaD}{3.2}
\newcommand{\EaPPlus}{16}
\newcommand{\EaTPlus}{28}
\newcommand{\EaBand}{1500}
\newcommand{\EaBoard}{%
\chessboard[setfen={5rk1/4qppp/4p3/3p3b/3P4/3QP3/Pr3PPP/R3R1K1 w - - 0 22}, boardfontsize=14pt, pgfstyle=straightmove, arrow=to, linewidth=0.07em, color=red!75!black, markmoves={e3-e4}, pgfstyle=straightmove, arrow=to, linewidth=0.07em, color=green!50!black, markmoves={a2-a4}]
}
\newcommand{\EbFen}{\fen{3q1rk1/2n2p1p/1p1p1npQ/1PpPp3/2P1P3/r1N2B1P/5PP1/2R2RK1 b - - 2 22}}
\newcommand{\EbMover}{Black}
\newcommand{\EbPhase}{middlegame}
\newcommand{\EbPlayed}{22\ldots{}\symrook{}b3}
\newcommand{\EbTop}{22\ldots{}\symknight{}ce8}
\newcommand{\EbPlayedUci}{\uci{a3b3}}
\newcommand{\EbTopUci}{\uci{c7e8}}
\newcommand{\EbLineP}{22\ldots{}\symrook{}b3 23.~\symqueen{}e3 \symrook{}a3 24.~\symrook{}a1}
\newcommand{\EbLineT}{22\ldots{}\symknight{}ce8 23.~\symqueen{}d2 \symqueen{}a8 24.~\symqueen{}b2}
\newcommand{\EbWpBefore}{54}
\newcommand{\EbWpP}{48}
\newcommand{\EbWpT}{49}
\newcommand{\EbD}{1.2}
\newcommand{\EbPPlus}{9}
\newcommand{\EbTPlus}{39}
\newcommand{\EbBand}{1500}
\newcommand{\EbBoard}{%
\chessboard[setfen={3q1rk1/2n2p1p/1p1p1npQ/1PpPp3/2P1P3/r1N2B1P/5PP1/2R2RK1 b - - 2 22}, boardfontsize=14pt, pgfstyle=straightmove, arrow=to, linewidth=0.07em, color=red!75!black, markmoves={a3-b3}, pgfstyle=straightmove, arrow=to, linewidth=0.07em, color=green!50!black, markmoves={c7-e8}]
}
\newcommand{\EaMedian}{3.1}
\newcommand{\EbMedian}{1.2}


\title{Where do his choices differ from players 200 points above him?\\[0.3em]
\large Twenty-two kinds of move in his 10-minute games, against Maia-2 at the band above,
corrected on Lichess and priced by Stockfish}
\author{Chess Coach research lab}
\date{}

\begin{document}
\maketitle

\begin{abstract}
Which kinds of move does he play more often, or less often, than players \Above{} points above
him would in his own positions, and which of those differences cost him? Over \Positions{} of his
moves in \Games{} 10-minute games, Maia-2 at the band above, corrected for its bias on
\LichessPositions{} positions of Lichess players, gives the expected count of each of \Kinds{}
kinds of move, and Stockfish prices each difference. By the rule fixed before any position was
scored, the study gives no verdict, because its known-fault control failed. He leaves a piece
hanging no more often than the band above would ($R$ \HangsRFirst{} and \HangsRSecond{} on the two
halves, both intervals holding 1), and less often than players at his own rating would
(\HangsRhisFirst{} and \HangsRhisSecond, both intervals below 1). When he does leave one and the
band above would not, it costs him \HangsRelFirst{} and \HangsRelSecond{} win\% more than his
other differences. The null control passed. Without a verdict, four differences show on both
halves: he plays pawn breaks about half again as often as the band above ($R$ \BreakRFirst{} and
\BreakRSecond), retreats less often (\RetreatRFirst{} and \RetreatRSecond), and gives check
(\CheckRFirst{} and \CheckRSecond) and doubles or isolates his own pawns (\WeakRFirst{} and
\WeakRSecond) more often. None of the four costs him more than his other differences. Added after
the run: among Lichess players, retreats rise and checks fall with rating while pawn breaks show
no trend, so his pawn breaks look like his own way of playing, and his fewer retreats and extra
checks like the level he is climbing from.
\end{abstract}

\section{The question}

A coach who has watched a player for a while can say what he does differently from stronger
players: he trades too readily, he rarely retreats, he pushes the pawns in front of his own king.
That kind of remark changes how a player plays, and it is not an engine's top move. This study
asks it of his games with the app's own instruments. Which kinds of move does he play more often,
or less often, than players \Above{} points above him would in the same positions, and which of
those differences cost him? A difference that shows in both halves of his games and costs him
more than his other differences is a habit worth coaching. One that shows but costs no more is
his style, worth knowing but not worth changing, and one that shows in only one half is not a
habit.

The comparison needs a stand-in for the stronger players in his own positions. Maia-2, the app's
model of human moves, was trained on Lichess games to predict the move a player at a rating band
makes. At the band above his it gives each legal move a probability, so it says how often players
\Above{} points above him would play each kind of move in exactly his positions. A model's
probabilities are not players' choices, so the study first measures, on Lichess games at those
bands, how often players play each kind against how often Maia-2 predicts, and corrects for the
difference. Stockfish then prices each difference.

\section{How it was measured}

\paragraph{His positions.} His \Games{} 10-minute games in the app's archive at commit
\texttt{\AppCommit}, in order of their end time: the first \GamesFirst{} (\FirstFrom{} to
\FirstTo) and the last \GamesSecond{} (\SecondFrom{} to \SecondTo) are the two halves. Each of his
moves enters when the position is past the opening by the app's divider, he is not in check and
has at least two legal moves, his win\% before the move is between \WinLow{} and \WinHigh{} by
the archive's evaluation (the game is still open), and his clock shows at least \ClockFloor{}
seconds. That leaves \Positions{} positions from \GamesWithPositions{} games, \PositionsFirst{} in
the first half and \PositionsSecond{} in the second. Most of the rest are opening moves
(\SkipOpening), moves in check (\SkipCheck) or moves in decided games (\SkipDecided).

\paragraph{The bands.} Maia-2, the app's model of human moves, gives every legal move a
probability at a rating band, conditioned on the opponent's band. His band is his rating at the
game with the app's offset of \Offset{} points to Lichess, and the band above is his rating plus
\Above{}, the same way. His band was \HisBandsFirst{} in the first half and \HisBandsSecond{} in
the second, and the band above \AboveBandsFirst{} and \AboveBandsSecond. Maia-2's lowest band
holds every rating below 1100, so in \OneStepGames{} of his earliest games the band above is one
band up rather than two.

\paragraph{The kinds of move.} Every legal move in every position is read by the app's code
(Table~\ref{tab:kinds}). Of the \Kinds{} kinds, \MoveKinds{} come from the move itself,
\IdeaKinds{} from the insights the app reads for a move as the board shows it, and one, a move
that leaves a piece standing to be taken, is the control. A move can be of several kinds. A pawn left to be taken
does not count as hanging, and a capture of a knight or more that nothing takes back counts as a
tactic even when it recaptures in a trade.

\KindTable

\paragraph{The rate.} For a kind and a half, $N$ is the number of his moves of that kind. Maia-2's
expected count sums, over his positions, the probability it gives at the band above to the moves
of that kind. Maia-2 is a model, so it may predict a kind more or less often than players at its
band play it. To correct for that, \LichessMonth's Lichess games at 10 minutes give the first
\PerBand{} player-games at each band from \BandLowest{} to \BandHighest{} (the lowest band counted
from \LichessFloor), and up to \PerPlayer{} positions from each, drawn among those that pass his
filters, the win\% filter by Stockfish at \FilterNodes{} nodes, with at most \Tries{} tried:
\LichessPositions{} positions from \LichessWithPositions{} player-games. At each band, $\kappa$ is
the number of moves of the kind those players played over the number Maia-2 at that band
expected. His expected count $E$ multiplies each position's predicted share by $\kappa$ at that
position's band above, and $R = N/E$ is how often he plays the kind relative to players
\Above{} points above him, in his positions.

\paragraph{The cost.} $T$ is the move Maia-2 at the band above rates most likely and $P$ his move.
Stockfish through the app's engine service, at \CostNodes{} nodes with its hash cleared and no
tablebase, scores each, and $D$ is his win\% after $T$ minus after $P$ (zero where they are the
same move). For a kind he plays too often, the instances are the positions where $P$ is of that
kind and $T$ is not; for one he plays too rarely, where $T$ is and $P$ is not. The relative cost
is the mean $D$ over the instances minus the mean $D$ over his other differences, the positions
where $P$ is not $T$ and neither is of the kind.

\paragraph{Intervals.} \Draws{} bootstrap draws resample his games within each half and the
Lichess player-games within each band; intervals are the 95\% percentile intervals.

\paragraph{The rule, fixed before any position was scored.} Two controls come first, and if
either fails on either half the study gives no verdict. The known fault: a move that leaves a
piece standing to be taken must be over-used on both halves, and cost more than his other
differences on both. The null: a player who plays a move drawn from Maia-2's own policy at the
band above, compared with Maia-2's expected count without $\kappa$, must show no kind over- or
under-used. Then a kind is \emph{over-used} when $R$'s interval lies above 1 on both halves and
\emph{under-used} when it lies below 1 on both. In that direction it \emph{costs} when the
relative cost's interval lies above 0 on both halves and is \emph{cheap} when it lies below 0 on
both. The verdict is a \emph{costly habit}, a \emph{habit that costs no more}, a \emph{cheap
habit}, or \emph{no habit}. A kind is \emph{not enough data} when $E$ is under \MinCount{} in
either half, or, for the cost, when it has fewer than \MinCount{} instances in either half, and
\emph{not measurable} when the correction averaged over a half's positions lies outside
\KappaLow{} to \KappaHigh.

\section{The controls}

\paragraph{The known fault fails.} He leaves a piece hanging, by the app's reading, in
\HisHangsFirst\% of his positions in the first half and \HisHangsSecond\% in the second:
\HangsNFirst{} and \HangsNSecond{} moves. Against the band above, $R$ is \HangsRFirst{}
(\HangsRFirstCi) and \HangsRSecond{} (\HangsRSecondCi). Both intervals hold 1, so the control fails
on both halves, and by the rule the study gives no verdict on any kind. The cost half of the
control holds: where he leaves a piece hanging and the band above's most likely move does not,
his move gives up \HangsOverDFirst{} and \HangsOverDSecond{} win\% on average, \HangsRelFirst{}
(\HangsRelFirstCi) and \HangsRelSecond{} (\HangsRelSecondCi) more than his other differences.

\paragraph{Why it fails.} The control assumed he leaves pieces hanging more often than stronger
players, and he does not. Against Maia-2 at his own band, corrected the same way, $R$ is
\HangsRhisFirst{} (\HangsRhisFirstCi) and \HangsRhisSecond{} (\HangsRhisSecondCi), both below 1:
he leaves a piece hanging less often than players at his rating would in his positions. The
reading itself follows rating. Added after the run, Lichess players leave a piece hanging in
\LiHangsLow\% of their positions at \BandLowest{} and \LiHangsHigh\% at \BandHighest{}
(Table~\ref{tab:trend}). The control also asked for more than this sample can show. In his
positions the expected count at his band is \StepFirst{} times the band above's in the first half
and \StepSecond{} times in the second, while $R$'s interval spans about $\pm$\SpreadFirst\% and
$\pm$\SpreadSecond\% of $R$. A player who left pieces hanging exactly as often as players at his
band would have shown $R$ near \StepFirst{} in the first half, and an interval as wide as his
around it would still hold 1.

\paragraph{The null passes.} A player who plays a move drawn from Maia-2's policy at the band
above in each of his positions shows no kind over- or under-used on both halves. On single
halves, \NullMisses{} of the \NullIntervals{} intervals miss 1, about what chance gives at 95\%.

\section{What the positions describe}

With no verdict, this section describes; nothing in it is a habit by the rule.

\VerdictTable

\paragraph{Four differences on both halves.} Four kinds lie on one side of 1 on both halves
(Table~\ref{tab:verdicts}). He plays pawn breaks about half again as often as the band above:
$R$ is \BreakRFirst{} (\BreakRFirstCi) and \BreakRSecond{} (\BreakRSecondCi). He retreats less
often (\RetreatRFirst{} and \RetreatRSecond), and he gives check (\CheckRFirst{} and \CheckRSecond)
and doubles or isolates his own pawns (\WeakRFirst{} and \WeakRSecond) more often. None of the
four costs him more than his other differences: each relative cost's interval holds 0 on both
halves. On average his extra checks give up \CheckDFirst{} and \CheckDSecond{} win\% against the
band above's move, against \CheckRefFirst{} and \CheckRefSecond{} for his other differences, though the
interval of the difference holds 0. Captures, pawn pushes, even trades and tactics, \OneHalf{}
more kinds, differ on one half only.

\paragraph{His level or his own.} Against Maia-2 at his own band (Table~\ref{tab:level}), pawn
breaks (\BreakRhisFirst{} and \BreakRhisSecond), his own pawn weaknesses (\WeakRhisFirst{} and
\WeakRhisSecond) and even trades (\TradeRhisFirst{} and \TradeRhisSecond) are more frequent on both
halves, so they set him apart from players at his rating as well. Retreats and checks do not:
against his own band each is within the noise on one half. Even trades stand out against the band
above only in the second half, where $R$ is \TradeRSecond{} (\TradeRSecondCi).

\LevelTable

\PhaseTable

\paragraph{By phase.} His fewer retreats are a middlegame difference: $R$ is \RetreatMid{}
(\RetreatMidCi) in the middlegame and \RetreatEnd{} (\RetreatEndCi) in the endgame. He also moves
his king less often in the middlegame, castling aside, at \KingMid{} (\KingMidCi). His pawn breaks
stand out in both phases, at \BreakMid{} and \BreakEnd{} (Table~\ref{tab:phase}).

\CostTable

\paragraph{The cost of each difference.} Table~\ref{tab:cost} gives the instances in each
direction. His move is the band above's most likely move in \TopAboveFirst\% and
\TopAboveSecond\% of his positions; where it is not, it gives up \DiffDFirst{} and \DiffDSecond{}
win\% on average. Where he plays a pawn break and the band above would not, his move gives up
\BreakDFirst{} and \BreakDSecond{} win\%, \BreakPerGameFirst{} and \BreakPerGameSecond{} times a
game; where the band above would retreat and he does not, \RetreatDFirst{} and \RetreatDSecond,
\RetreatPerGameFirst{} and \RetreatPerGameSecond{} times a game. Neither is far from his other
differences.

\KappaTable

\paragraph{Maia-2's bias.} Maia-2 is close to the players at its bands. Where a band had enough
moves of a kind, $\kappa$ lies between \KappaMin{} and \KappaMax{} (Table~\ref{tab:kappa}); in
\KappaPooled{} of the \KappaCells{} cells the pooled bias stands in. Averaged over a half, the
correction lies between \CorrMin{} and \CorrMax, so every kind is measurable.

\section{Added after the run}

These descriptions were added once the rule had been applied; none was named before the run.

\TrendTable

\paragraph{The bands on Lichess.} Among Lichess players (Table~\ref{tab:trend}), retreats rise
from \LiRetreatLow\% of positions at \BandLowest{} to \LiRetreatHigh\% at \BandHighest, and checks
fall from \LiCheckLow\% to \LiCheckHigh\%: both move with rating, the way his differences point.
Pawn breaks show no trend, between \LiBreakMin\% and \LiBreakMax\% at every band, while he plays
one in \HisBreakFirst\% and \HisBreakSecond\% of his positions, where Maia-2 at the band above
expects \MaiaBreakFirst\% and \MaiaBreakSecond\%. His pawn breaks look like his own way of
playing, and his fewer retreats and extra checks like the level he is climbing from.

\paragraph{Whose moves his are.} Maia-2 at his own band gives his moves a little more probability
than at the band above, on both halves: \LikeFirst\% and \LikeSecond\% less a move at the band
above, both intervals clear of zero. His move is his own band's most likely move \TopHisFirst\%
and \TopHisSecond\% of the time, and the band above's \TopAboveFirst\% and \TopAboveSecond\%. He
plays like his band, with fewer pieces left hanging.

\section{Two positions}

Two of his positions show a pawn break the band above would not play and a retreat it would
(Figure~\ref{fig:examples}). They illustrate; they test nothing.

\begin{figure}[ht]
\centering
\begin{subfigure}[t]{0.48\textwidth}
\centering
\EaBoard\\[0.4em]
{\footnotesize FEN \EaFen}
\caption{He is \EaMover{}, in the \EaPhase{} at \EaWpBefore\% to win. He played the pawn break
\EaPlayed{} (red), which Maia-2 at \EaBand{} gives \EaPPlus\%; it rates \EaTop{} (green) most
likely, at \EaTPlus\%. Stockfish: after \EaLineP{} he has \EaWpP\%, after \EaLineT{} \EaWpT\%.
$D$ is \EaD, near the median of these instances (\EaMedian).}
\label{fig:break}
\end{subfigure}\hfill
\begin{subfigure}[t]{0.48\textwidth}
\centering
\EbBoard\\[0.4em]
{\footnotesize FEN \EbFen}
\caption{He is \EbMover{}, in the \EbPhase{} at \EbWpBefore\% to win. Maia-2 at \EbBand{} rates
the retreat \EbTop{} (green) most likely, at \EbTPlus\%; he played \EbPlayed{} (red), which it
gives \EbPPlus\%. Stockfish: after \EbLineP{} he has \EbWpP\%, after \EbLineT{} \EbWpT\%. $D$ is
\EbD, near the median of these instances (\EbMedian).}
\label{fig:retreat}
\end{subfigure}
\caption{His second-half positions. Of the middlegame instances where neither move captures or
gives check, the one whose $D$ is nearest the median $D$ of all the second half's instances: left,
a pawn break he played and the band above would not; right, a retreat the band above would play
and he did not. Each line is the engine's, at \CostNodes{} nodes.}
\label{fig:examples}
\end{figure}

\section{What the app can take}

\begin{itemize}[nosep]
\item \textbf{Nothing yet.} With the known fault failing, no difference here is a habit by the
  rule, so the coach does not write to any of them.
\item \textbf{A control he has.} A next round needs a positive control that shows the method
  finds a difference of the size it reports: a fault he is known to have against stronger
  players, or a planted one, a player made to play a kind more often by a set factor. Its rule is
  fixed first.
\item \textbf{The candidates.} Pawn breaks, the clearest difference on both halves and against
  his own band, and retreats in the middlegame, the difference that moves with rating on Lichess.
  Neither costs him more than his other differences here, so either would be his style rather
  than a fault unless a next round says otherwise.
\item \textbf{A model of him.} A network trained on his own games, the plan's next step, needs a
  GPU.
\end{itemize}

\section{Caveats}

\begin{itemize}[nosep]
\item \textbf{Maia-2 stands in for the players above him.} $\kappa$ corrects its average bias for
  each kind of move, measured in Lichess players' positions, not in his.
\item \textbf{The most likely move stands for what the band above plays} when the cost is
  measured.
\item \textbf{The halves are at different ratings.} He climbed through the archive, so a
  difference that shows in both halves held as he improved; one that shows in only one may have
  come or gone.
\item \textbf{A kind of move is a description, not a mistake.} Playing it more often than the
  band above is a difference; the cost says whether it matters.
\item \textbf{The app's readings at the pinned commit.} The kinds are the app's readings as the
  board shows them, without an engine. A piece the reading calls hanging may be defended by a
  tactic, and a recapture can count as a tactic.
\item \textbf{Added after the run.} The Lichess rates by band, the likelihood of his moves at each
  band, the size of the step between bands and the two positions were looked at once the rule had
  been applied. They describe; a rule built on them needs games this study did not see.
\end{itemize}

\end{document}
